\begin{longtable}[c]{|m{5cm}|m{10cm}|}
\caption{Configuration Parameters of the Operator Submodule}
\label{tab:mcpwm-param-operator}
    \\\hline
    \rowcolor{lightgray}
   \textbf{Submodule} & \textbf{Configuration Parameter or Option} \\ \hline 

\endfirsthead\hline
 \rowcolor{lightgray}
   \textbf{Submodule} & \textbf{Configuration Parameter or Option}\\
 \hline 
\endhead
\endfoot
    PWM Generator &\vspace{-1em}
    \begin{itemize}
        \item Set up the PWM duty cycle for PWM\textcolor{red}{\it{x}}A and/or PWM\textcolor{red}{\it{x}}B output.
        \item Set up at which time the timing events occur.
        \item Define what action should be taken on timing events:
        \begin{itemize}
           \item Switch high or low of PWM\textcolor{red}{\it{x}}A and/or PWM\textcolor{red}{\it{x}}B outputs
           \item Toggle PWM\textcolor{red}{\it{x}}A and/or PWM\textcolor{red}{\it{x}}B outputs
           \item Take no action on outputs
        \end{itemize}
        \item Use direct s/w control to force the state of PWM outputs
        \item Add a dead time to raising edge and/or failing edge on PWM outputs.
        \item Configure update method for this submodule.\vspace{-1em}
    \end{itemize}\\\hline
    
    Dead Time Generator &\vspace{-1em}
    \begin{itemize}
        \item Control of complementary dead time relationship between upper and lower switches.
        \item Specify the dead time on rising edge.
        \item Specify the dead time on falling edge.
        \item Bypass the dead time generator module. The PWM waveform will pass through without inserting dead time.
        \item Allow PWM\textcolor{red}{\it{x}}B phase shifting with respect to the PWM\textcolor{red}{\it{x}}A output.
        \item Configure updating method for this submodule.\vspace{-1em}
    \end{itemize}\\\hline
    
    PWM Carrier &\vspace{-1em}
    \begin{itemize}
        \item Enable carrier and set up carrier frequency.
        \item Configure duration of the first pulse in the carrier waveform.
        \item Set up the duty cycle of the following pulses.
        \item Bypass the PWM carrier module. The PWM waveform will be passed through without modification.\vspace{-1em}
    \end{itemize}\\\hline
    
    Fault Handler &
    \begin{itemize} \vspace{-1em}   
        \item Configure if and how the PWM module should react the fault event signals.
        \item Specify the action taken when a fault event occurs:
        \begin{itemize}
            \item Force PWM\textcolor{red}{\it{x}}A and/or PWM\textcolor{red}{\it{x}}B high.
            \item Force PWM\textcolor{red}{\it{x}}A and/or PWM\textcolor{red}{\it{x}}B low.
            \item Configure PWM\textcolor{red}{\it{x}}A and/or PWM\textcolor{red}{\it{x}}B to ignore any fault event.
        \end{itemize}
        \item Configure how often the PWM should react to fault events:
        \begin{itemize}
            \item One-shot
            \item Cycle-by-cycle
        \end{itemize}
        \item Generate interrupts.
        \item Bypass the fault handler submodule entirely.
        \item Set up an option for cycle-by-cycle actions clearing.
        \item If desired, independently-configured actions can be taken when time-base counter is counting down or up.\vspace{-1em}

    \end{itemize}\\\hline
    
\end{longtable}
